% Source: https://artofproblemsolving.com/community/c3h1974938_large_board
% Stable identity: aops-topic-1974938-q01
A very large board is made of unit squares ($1$ cm by $1$ cm). The positive integers are written in the squares in increasing order, starting with $1$ and forming square ``layers'' around $1$: the first layer is the square containing $1$; the second layer consists of the $8$ squares around it, containing $2$ to $9$; and the third layer starts with $10$ in the square to the left of $9$, and so on, each layer running once around the previous one in the same direction. The figure shows the start of the board (in the original, $2$ is to the left of $1$ and the numbers run counterclockwise).

\begin{asy}
usepackage("amsmath");
size(4.5cm);
defaultpen(fontsize(10pt));
int[][] g={{0,0,1},{-1,0,2},{-1,-1,3},{0,-1,4},{1,-1,5},{1,0,6},{1,1,7},{0,1,8},{-1,1,9},{-2,1,10},{-2,0,11},{-2,-1,12}};
for(int x=-2;x<=2;++x) for(int y=-2;y<=2;++y) draw(shift(x-0.5,y-0.5)*unitsquare,gray(0.6));
for(int[] t:g) label("$"+string(t[2])+"$",(t[0]-0.22,t[1]+0.22),fontsize(9pt));
draw((0,0)--(-2,0)--(-2,-1)--cycle,linewidth(0.9)+blue);
dot((0,0),blue+3pt);dot((-2,0),blue+3pt);dot((-2,-1),blue+3pt);
label("\mbox{}",(min(currentpicture,user=true)+max(currentpicture,user=true))/2,black); // keeps the text after the figure black when LaTeX typesets the labels inline
\end{asy}

To find the distance from square $1$ to a square $N$ not in its row or column, we use the right triangle with one horizontal side, one vertical side, and hypotenuse joining the centre of square $1$ to the centre of square $N$. For square $12$ this triangle has vertices at the centres of squares $1$, $11$ and $12$ (shown); its horizontal side is $2$ cm and its vertical side is $1$ cm.
\begin{enumerate}
\item[(a)] What is the last number of the third layer?
\item[(b)] Find the horizontal and vertical sides of the triangle for square $33$.
\item[(c)] Find the horizontal and vertical sides of the triangle for square $2019$.
\end{enumerate}
